\documentclass[11pt, a4paper]{article}
% ---- Core Geometry & Margins ----
\usepackage[margin=1in, top=1.2in, bottom=1.2in]{geometry}

% ---- Font & Micro-Typography ----
\usepackage[T1]{fontenc}
\usepackage{lmodern}
\usepackage{microtype}

% ---- Language & Math ----
\usepackage[english]{babel}
\usepackage{amsmath, amssymb, amsthm}
\usepackage{csquotes}

% ---- Formatting & Aesthetics ----
\usepackage{booktabs}
\usepackage{graphicx}
\usepackage{xcolor}

% ---- Color Palette ----
\definecolor{navyblue}{RGB}{16, 55, 87}
\definecolor{darkslate}{RGB}{53, 64, 83}

% ---- Bibliography ----
\usepackage[backend=biber, style=numeric-comp, sorting=nyt, giveninits=true]{biblatex}
\addbibresource{references.bib}

% ---- Hyperlinks (load LAST) ----
\usepackage[
  pdfencoding=auto,
  colorlinks=true,
  linkcolor=navyblue,
  citecolor=navyblue,
  urlcolor=navyblue
]{hyperref}

% ---- Metadata ----
\title{Skills and Competences for the Sustainability Engineer}
\author{Tadouanla Guetchuin Billy \\ 28177}

\date{\today}

% ====================================================================================
\usepackage{parskip}
\begin{document}
\maketitle
\pagenumbering{roman}
\setcounter{page}{1}

\tableofcontents
\newpage
\pagenumbering{arabic}
\setcounter{page}{1}
\begin{abstract}
  Today's engineering education is pressurised to respond to the Sustainable Development Goals (SDGs),
  yet the literature lacks consensus on which competences future engineers actually need. This short
  paper analyses\cite{beagon2023preparing} European focus-group study---covering 86 students,
  academics, and employers across Denmark, Finland, France, and Ireland---to examine three competences
  relevant to a sustainability engineer:\textbf{systems thinking}, \textbf{life cycle thinking}, and
  \textbf{interpersonal competence}. Upon analysis, the paper shows strong stakeholder agreement that
  systems thinking and interpersonal competence are essential, while life cycle thinking is recognised
  primarily by students  and remains implicit. The analysis questions the notion of a
  \textquote{typical}
  sustainability engineer, arguing that no single profile exists; instead, the typical engineer is a
  technically grounded professional who adds \textbf{holistic} and  \textbf{human-centred} competences
  to disciplinary depth. The analysis concludes that a gap is not in identifying competences but in
  integrating them into overcrowded curricula.
\end{abstract}

\section{Introduction}
Engineering problems are described in the literature as \textquote{wicked} problems---complex,
multifaceted, and resistant to purely technical solutions\cite{beagon2023preparing}. Future Engineers
are expected to act as catalyst of technical innovation and lead on socioeconomic issues. The UN
Sustainable Development Goals (SDGs) provide a framework for this shift, covering 169 targets across
environmental, social, and economic domains. However, no agreement on which competences engineering
graduate should develop to meet these challenges.\\

This paper analyses selected literature on the skills and competences required by engineers
for sustainability, using  \textbf{Beagon at al. (2023)} as primary source. Three competences for the
typical engineer for sustainability is addressed. The paper first outlines the methodological approach,
then analyses systems thinking, life cycle thinking, and interpersonal competence, before discussing
the concept of the \textquote{typical} engineer.

\section{Methodology}
The study employed a qualitative design using twelve focus groups conducted in Denmark,
Finland, France, and Ireland between March and July 2019, involving 86 participants from three
stakeholder groups: engineering academics, employers, and students. Data was analysed through an open
coding process, producing 54 competences grouped into six categories: Fundamental Technical Skills,
Application Skill, Outward facing People-Oriented Skills, Inward-facing Ways of thinking, World view,
and Character and Ethical Orientation and mapped against the UNESCO (2017) key competences framework.
\\

For the this paper, three competences were selected from the study's findings to represent
different sustainability practice: system thinking (a way of thinking), life cycle thinking (an
application skill), and interpersonal competence (a people-oriented skill).

\newpage
\section{Analysis of Three Key Competences}
\subsection{System Thinking}
It appears in the findings under the category of \textquote{Inward-Facing Ways of Thinking,}
operationalised as analytical and holistic thinking. Holistic thinking is defined by participants as
\textquote{the capacity to make a good synthesis and see the whole picture.} All three stakeholder
groups---students, academics, and employers---identified systems thinking as a key competence.\\

Its relevance to sustainability is grounded in the nature of sustainability problems
themselves. Because engineering challenges are \textquote{ill-defined, open-ended, cross-disciplinary,
  complex socio-technical problems,} an engineer who optimises a single component in isolation risks
shifting burden elsewhere in the system One academic quoted in the study described the shift required:
students in university \textquote{do what we say to do} without a risk or responsibility, but in
professional practice \textquote{they have a responsibility and they have to be able to justify their
  actions and apply global thinking}. Another stakeholder emphasised the ability to
\textquote{perceive the worldwide situation, not only from the engineer's perspective, but also from
  other viewpoints.}\\

System thinking therefore functions as the cognitive foundation of sustainability engineering: it is the competence that allows an engineer to situate a technical decision within environmental, social, and economic systems simultaneously.

\subsection{Life Cycle Thinking}
This is the ability to consider a product, process, or infrastructure across its entire life---from
material extraction through use to disposal.\\

Its relevance is closely tied to what the study identifies as a notable gap:
\textbf{anticipatory competence}. Beagon et al.found that none of the stakeholder groups explicitly identified anticipatory competence---the ability to envision future scenarios and estimate
consequences---despite it begin one of UNESCO' (2017) eight key competences. The authors argue that
this competence is \textquote{not entirely absent} because similar ideas appear implicitly in other
competences, specifically naming life cycle thinking, lifelong learning, and time management as
carriers of future-oriented thinking.

A striking finding is the divergence between stakeholder groups. Life cycle thinking was identified by
\textbf{students} as a necessary competence, but not explicitly by \textbf{academics} or
\textbf{employers}. This suggests that students---who will inherit the long-term consequence of today's
design decisions---may already be internalising the future-oriented perspective that the literature
consider essential, even where industry and curricular have not yet formalised it.

\subsection{Interpersonal Competence}
This competence mainly encompasses communication, collaboration, teamwork, and related abilities such
as listening, emotional intelligence, and empathy. The study describes interpersonal competence as
\textquote{competences associated with engaging with people, on a general level, regardless of their
  discipline}.


\newpage
This is arguably the study's most emphatic finding. Its relevance to sustainability is practical:
sustainability solutions require negotiation across diverse value sets. As one employer explained,
\textquote{engineers will also need to manage divers project teams, there are no longer only engineers on teams[..] They now include environmentalists, economists, institutional capacity builders, all kinds of people.}

Yet the study also documents a persistent gap. One employer stated that \textquote{some of them come
  in lacking in those skills, particularly compared to business graduates.}Academics noted that while
collaboration is taught implicitly through project work, competences such as empathy and emotional
intelligence are \textquote{much needed but not sufficiently promoted in engineering curricula.}
The study concludes that there needs to be \textquote{more emphasis and explicit assessment of
  interpersonal competences \dots{} rather than relying on gaining these competences implicitly through
  project work.}

\section{Discussion}
The term \textquote{typical} is difficult to apply to the sustainability engineer, and the literature
suggests it may be misleading if understood as a distinct job title. Beagon et al.'s findings support
three observations.

First, technical depth remains non-negotiable. One academic stated that \textquote{hard-core, solid
  technical skills are required—however, if engineers have technical skills only that is clearly
  inadequate,} while an employer warned that universities must not compromise technical training. The
\textquote{typical} sustainability engineer is therefore not a generalist who has traded depth for
breadth; the study explicitly rejects this, noting that \textquote{not all people can be generalists}
and that each engineer \textquote{has to understand how his/her special expertise is connected to
  others.}

Second, the competences identified are additive rather than substitutive. Systems thinking, life
cycle thinking, and interpersonal competence are layered on top of disciplinary expertise. This
aligns with the T-shaped profile: depth in one discipline (the vertical bar) combined with the
ability to connect across disciplines and stakeholders (the horizontal bar). The \textquote{typical}
sustainability engineer is thus a mechanical, civil, or electrical engineer who applies
sustainability competences within their domain—not a separate profession.

Third, the study reveals that the \textquote{typical} profile is still contested. Employers called
for mathematics and interpretation skills and for graduates who \textquote{challenge the status quo};
academics emphasised conceptual understanding and negotiation; students valued organisation,
leadership, and life cycle thinking that neither of the other groups named. The absence of
anticipatory competence from all groups further complicates the picture, since a future-oriented
competence is precisely what sustainability demands.

The most defensible definition, therefore, is that the \textquote{typical} sustainability engineer is
defined less by a fixed role than by a mindset: a technically competent professional who refuses
siloed thinking, considers long-term consequences, and can work with people who do not share their
disciplinary background.

\newpage
\section{Conclusion}
This paper analysed Beagon et al.'s (2023) study of engineering competences for the SDGs and discussed
three relevant competences. Systems thinking provides the cognitive framework for complex
socio-technical problems, life cycle thinking extends that framework across time, and interpersonal
competence enables engineers to negotiate sustainability solutions with diverse stakeholders. The
\textquote{typical} sustainability engineer is therefore not a separate profession, but a technically
grounded engineer who adds holistic, future-oriented, and human-centred competences to disciplinary
depth\cite{beagon2023preparing}. The real barrier is not identifying these competences, but
integrating them into already overcrowded curricula.

\newpage
\printbibliography[
  heading=bibintoc,
  title={References}
]
\end{document}
